% Ch15 self-study -- "I do / You do" ladder for Zumdahl chapter 15.
% FOURTEEN ladders (56 questions) plus a ten-question mixed set = 66
% questions, fourteen figures. Completes the self-study series.
%
% SCOPE, read straight off the CED PDF on 2026-08-23. This chapter
% carries the most intricate scope in the whole course, so it is worth
% recording precisely:
%
%   ASSESSED (8.4.A.1-A.4): pH of any mixture of strong/weak acid and
%   base -- excess-reagent case, buffer case (via H-H), and the
%   equimolar/equivalence case (via the conjugate's Kb or Ka).
%   ASSESSED (8.5): titration curves, equivalence point, indicators.
%   ASSESSED (8.8, 8.10): why buffers work, and buffer capacity --
%   both argumentation-flavoured, not computational.
%   ASSESSED (8.9): the Henderson-Hasselbalch equation itself.
%
%   EXCLUDED (8.9): "Computation of the change in pH resulting from the
%   addition of an acid or a base to a buffer" AND "Derivation of the
%   Henderson-Hasselbalch equation".
%   EXCLUDED (8.5): "Computation of the concentration of each species
%   present in the titration curve for POLYPROTIC acids" -- the CED
%   states explicitly that the same computations FOR MONOPROTIC ACIDS
%   are in scope.
%
% The 8.9 exclusion is narrow and easy to over-read: mixing a weak acid
% with a strong base and finding the resulting pH is 8.4.A.2 and IS
% assessed. What is excluded is being handed an existing buffer and
% asked for the pH CHANGE on adding something. Ladder 6 says this
% explicitly, because the two calculations are nearly identical and a
% student who over-reads the exclusion would skip required material.
\documentclass{apchem-worksheet}

\apcunitword{Chapter}
\apcunit{15}
\apcunittitle{Acid--Base Equilibria}
\apcdoctitle{Self-Study \textbullet\ Chapter 15, I~do / You~do}
\apcsubtitle{Zumdahl \S15.1--15.6 \textbullet\ PDF pp.\ 756--792}

\begin{document}
\apctitleblock

\begin{apcnote}[How to use these notes --- read this first]
Each skill is a \textbf{ladder}: a fully worked example in the
solid-framed box, then \textbf{four} YOUR TURN questions in the dashed
box --- same skill, new numbers, no help. Work all four before checking
the gray \emph{check:} line; a tracker tick needs all four right first
time. A ten-question \textbf{mixed set} follows Ladder 14.

\textbf{This is the second half of Unit 8}, which at 11--15\% is the
exam's second-heaviest unit. Chapter 14 gave you weak acids in
isolation; this chapter puts them in mixtures --- which is where nearly
all the exam's Unit 8 marks actually live.

\textbf{What you must bring.} Chapter 14, Ladders 9, 12 and 16. Every
calculation here is a weak-acid equilibrium with an extra stoichiometry
step in front of it. If $K_{a}K_{b} = K_{w}$ and $\text{p}K_{a}$ are not
automatic, fix that first.

\textbf{The one structural idea.} Almost every problem in this chapter
is \textbf{two steps}: first a \emph{stoichiometry} step in moles ---
the strong reagent is consumed completely --- and then an
\emph{equilibrium} step on whatever is left. Students who try to do both
at once get lost; students who do them in order rarely do.

\textbf{A scope note that matters.} The CED excludes two specific
things here (Ladders 6 and 14 flag them precisely). Both exclusions are
narrower than they first appear, and over-reading either one would make
you skip required material --- so read those notes rather than
guessing.
\end{apcnote}

% =====================================================================
\section*{\sffamily Ladder 1 \textbullet\ The common-ion effect}
\zsec{15.1}

Add acetate to acetic acid and the acid ionises \emph{less}. That is
Le Ch\^{a}telier applied to a weak acid --- and it is the mechanism
behind every buffer in this chapter.

\begin{center}
\begin{tikzpicture}[font=\sffamily\footnotesize]
  % pure weak acid -- fill before stroke
  \fill[apcgray!22] (0.7,0.5) rectangle (3.3,2.5);
  \draw[very thick] (0.7,3.0) -- (0.7,0.5) -- (3.3,0.5) -- (3.3,3.0);
  \foreach \p in {(1.1,0.95),(1.9,1.5),(2.7,1.0),(1.5,2.1),(2.4,2.0)} {
    \node[font=\tiny] at \p {\ce{HA}};}
  \node[font=\tiny] at (2.0,0.75) {\ce{H+}};
  \node[align=center] at (2.0,-0.2)
    {pure weak acid\\a little ionises};

  \draw[-{Stealth[length=3mm]},very thick] (3.8,1.5) -- (5.1,1.5);
  \node[above,align=center,font=\sffamily\footnotesize] at (4.45,1.65)
    {add\\\ce{A-}};

  \begin{scope}[xshift=5.6cm]
    \fill[apcgray!22] (0.7,0.5) rectangle (3.3,2.5);
    \draw[very thick] (0.7,3.0) -- (0.7,0.5) -- (3.3,0.5) -- (3.3,3.0);
    \foreach \p in {(1.0,0.9),(1.7,1.4),(2.5,0.95),(1.3,2.0),
                    (2.2,1.9),(2.9,1.5)} {
      \node[font=\tiny] at \p {\ce{HA}};}
    \foreach \p in {(1.1,1.7),(2.0,2.3),(2.7,2.2)} {
      \node[font=\tiny] at \p {\ce{A-}};}
    \node[align=center] at (2.0,-0.2)
      {with added \ce{A-}\\ionisation \textbf{suppressed}};
  \end{scope}
  \node[font=\sffamily\footnotesize] at (5.6,-1.1)
    {$\ce{HA <=> H+ + A-}$ --- adding a product pushes the equilibrium
     \textbf{left}};
\end{tikzpicture}
\end{center}

\begin{workedexample}[I do: the same acid, two solutions]
\textbf{Compare \SI{0.10}{\Molar} acetic acid alone with
\SI{0.10}{\Molar} acetic acid that also contains \SI{0.10}{\Molar}
sodium acetate. $K_{a} = \num{1.8e-5}$.}

\textbf{Alone} (Chapter 14, Ladder 9): $x = \sqrt{K_{a}c} =
\num{1.34e-3}$, so pH $= 2.87$.

\textbf{With acetate present}, the ICE table starts differently ---
acetate is no longer zero:
\begin{center}
\renewcommand{\arraystretch}{1.15}
\begin{tabular}{@{}lccc@{}}
\toprule
& $[\ce{HA}]$ & $[\ce{H+}]$ & $[\ce{A-}]$ \\
\midrule
\textbf{I} & 0.10 & 0 & \textbf{0.10} \\
\textbf{C} & $-x$ & $+x$ & $+x$ \\
\textbf{E} & $0.10-x$ & $x$ & $0.10+x$ \\
\bottomrule
\end{tabular}
\end{center}

\[ K_{a} = \frac{x(0.10+x)}{0.10-x}
   \;\approx\; \frac{x(0.10)}{0.10} = x \]
So $x = K_{a} = \num{1.8e-5}$ and pH $= 4.74$.

\textbf{Nearly two pH units less acidic}, from adding a salt that
contains no \ce{H+} at all. The added acetate pushed the ionisation
equilibrium left, exactly as Chapter 13's Le Ch\^{a}telier predicts.

\textbf{Notice what the algebra collapsed to.} With comparable amounts
of \ce{HA} and \ce{A-}, the two concentrations nearly cancel and
$[\ce{H+}] \approx K_{a}$. That is the Henderson--Hasselbalch equation
in embryo (Ladder 3) --- and it is why buffers are so easy to compute
once you see the pattern.
\end{workedexample}

\begin{yourturn}[YOUR TURN 1 --- four questions]
\begin{enumerate}[leftmargin=*,itemsep=0.5em]
  \item Does adding \ce{NaF} to \ce{HF} solution raise or lower the pH?
        Explain. \workspace[1.8cm]
  \item Does it raise or lower the percent ionisation of \ce{HF}?
        \workspace[1.6cm]
  \item Does it change $K_{a}$? \workspace[1.4cm]
  \item Which ion in \ce{NaF} is the ``common'' ion, and why is the
        other irrelevant? \workspace[1.8cm]
\end{enumerate}
\selfcheck{(a) raises it --- equilibrium shifts left \quad (b) lowers
it \quad (c) no \quad (d) \ce{F-}; \ce{Na+} is a spectator}
\keyonly{\begin{apcanswer}(a) \textbf{Raises} the pH. \ce{F-} is a
product of $\ce{HF <=> H+ + F-}$, so adding it shifts the equilibrium
\textbf{left}, reducing $[\ce{H+}]$. (b) \textbf{Lowers} it --- less of
the \ce{HF} comes apart. (c) \textbf{No.} $K_{a}$ depends only on
temperature; the common ion moves the position of the equilibrium, not
the constant. (d) \textbf{\ce{F-}} --- it is common to both the salt
and the acid's equilibrium. \ce{Na+} is a spectator: it is the
conjugate of the strong base \ce{NaOH} and does not react with water
(Chapter 14, Ladder 14).\end{apcanswer}}
\end{yourturn}

% =====================================================================
\section*{\sffamily Ladder 2 \textbullet\ What a buffer is, and why it
works}
\zsec{15.2}

A buffer contains \textbf{appreciable amounts of both members of a
conjugate pair}. That is the whole definition, and everything else
follows from it.

\begin{center}
\begin{tikzpicture}[font=\sffamily\footnotesize]
  \node[font=\sffamily\small\bfseries] at (5.5,3.5)
    {a buffer has two jobs, and one reagent for each};
  % add acid
  \node[anchor=west,font=\sffamily\bfseries] at (0.5,2.8)
    {add \ce{H+}?};
  \node[anchor=west] at (2.6,2.8)
    {the \textbf{conjugate base} mops it up: \;
     $\ce{A- + H+ -> HA}$};
  % add base
  \node[anchor=west,font=\sffamily\bfseries] at (0.5,2.2)
    {add \ce{OH-}?};
  \node[anchor=west] at (2.6,2.2)
    {the \textbf{conjugate acid} mops it up: \;
     $\ce{HA + OH- -> A- + H2O}$};
  \draw[apcgray] (0.4,1.85) -- (11.0,1.85);
  \node[anchor=west] at (0.5,1.4)
    {either way the strong species is converted into a \emph{weak} one,
     so the pH barely moves};
  \node[anchor=west] at (0.5,0.85)
    {the \textbf{ratio} $[\ce{A-}]/[\ce{HA}]$ shifts a little --- and
     pH depends on the \emph{logarithm} of that ratio};
  \draw[apcgray] (0.4,0.5) -- (11.0,0.5);
  \node[anchor=west,font=\sffamily\itshape] at (0.5,0.05)
    {a buffer does not stop pH changing --- it makes the change small};
\end{tikzpicture}
\end{center}

\begin{workedexample}[I do: why the pH barely moves]
\textbf{This is CED topic 8.8, and it is an \emph{argumentation} topic
--- the marks are for the explanation, not a number.}

\textbf{The mechanism, stated properly.} A buffer solution contains a
large concentration of both members of a conjugate acid--base pair. The
conjugate base reacts with added acid; the conjugate acid reacts with
added base. Both reactions convert a \emph{strong} species into a
\emph{weak} one, and a weak species contributes far less to the pH.

\textbf{Why the logarithm matters.} pH tracks the \emph{ratio}
$[\ce{A-}]/[\ce{HA}]$, and only through a logarithm. Take a buffer that
is \SI{0.50}{\Molar} in each. Convert a tenth of the acid to base and
the ratio goes from $1.00$ to $0.60/0.40 = 1.5$ --- and
$\log 1.5 = 0.18$, so the pH moves by less than two-tenths of a unit.
The same amount of acid added to pure water would move the pH by
several units.

\textbf{What a buffer is not.} It is not a substance that ``absorbs''
\ce{H+} indefinitely, and it does not hold pH \emph{constant}. It
resists change until one component is used up --- which is Ladder 5.

\textbf{The exam sentence.} ``The solution contains appreciable amounts
of both \ce{HA} and \ce{A-}. Added \ce{H+} reacts with \ce{A-} to form
\ce{HA}, and added \ce{OH-} reacts with \ce{HA} to form \ce{A-}, so the
ratio changes only slightly and the pH is stabilised.'' Both directions,
or you have answered half the question.

\textbf{And what makes a mixture a buffer.} Not any acid and base ---
a \textbf{conjugate pair}. \ce{HCl} and \ce{NaOH} together are not a
buffer; they neutralise each other completely. \ce{CH3COOH} with
\ce{NaCH3COO} is, and so is \ce{NH3} with \ce{NH4Cl}.
\end{workedexample}

\begin{yourturn}[YOUR TURN 2 --- four questions]
\begin{enumerate}[leftmargin=*,itemsep=0.5em]
  \item Which of these is a buffer? \ce{HCl}/\ce{NaCl},
        \ce{HF}/\ce{NaF}, \ce{HCl}/\ce{NaOH}. \workspace[1.8cm]
  \item Which component reacts with added \ce{OH-}?
        \workspace[1.4cm]
  \item Write the equation for a buffer absorbing added \ce{H+}.
        \workspace[1.6cm]
  \item Does a buffer keep the pH exactly constant? Explain.
        \workspace[1.8cm]
\end{enumerate}
\selfcheck{(a) \ce{HF}/\ce{NaF} \quad (b) the weak acid \ce{HA} \quad
(c) $\ce{A- + H+ -> HA}$ \quad (d) no --- it makes the change small}
\keyonly{\begin{apcanswer}(a) \textbf{\ce{HF}/\ce{NaF}} --- the only
\textbf{conjugate pair}. \ce{HCl}/\ce{NaCl} is a strong acid with its
useless conjugate; \ce{HCl}/\ce{NaOH} simply neutralise. (b) The
\textbf{weak acid, \ce{HA}}: $\ce{HA + OH- -> A- + H2O}$. (c)
$\ce{A- + H+ -> HA}$. (d) \textbf{No.} It \emph{resists} change ---
added strong acid or base is converted to a weak species, so the ratio
$[\ce{A-}]/[\ce{HA}]$ shifts only slightly and, because pH depends on
the logarithm of that ratio, the pH moves very
little.\end{apcanswer}}
\end{yourturn}

% =====================================================================
\section*{\sffamily Ladder 3 \textbullet\ Henderson--Hasselbalch}
\zsec{15.2}

\[ \text{pH} = \text{p}K_{a}
   + \log\!\left(\frac{[\ce{A-}]}{[\ce{HA}]}\right) \]

\begin{center}
\begin{tikzpicture}[font=\sffamily\footnotesize]
  \node[font=\sffamily\small\bfseries] at (5.5,3.2)
    {read the ratio, not the concentrations};
  \foreach \y/\rat/\lg/\ph in {%
      2.55/{$[\ce{A-}] = [\ce{HA}]$}/{$\log 1 = 0$}/{pH $=$
        p$K_{a}$},
      2.0/{$[\ce{A-}] = 10[\ce{HA}]$}/{$\log 10 = 1$}/{pH $=$
        p$K_{a} + 1$},
      1.45/{$[\ce{A-}] = 0.1[\ce{HA}]$}/{$\log 0.1 = -1$}/{pH $=$
        p$K_{a} - 1$}} {
    \node[anchor=west] at (0.6,\y) {\rat};
    \node[anchor=west] at (4.2,\y) {\lg};
    \node[anchor=west] at (7.2,\y) {\ph};}
  \draw[apcgray] (0.5,2.9) -- (11.0,2.9);
  \draw[apcgray] (0.5,1.1) -- (11.0,1.1);
  \node[anchor=west] at (0.6,0.65)
    {because it is a \textbf{ratio}, you may use \textbf{moles}
     instead of concentrations --- the volume cancels};
  \node[anchor=west,font=\sffamily\itshape] at (0.6,0.1)
    {which is what makes the titration ladders (7--10) quick};
\end{tikzpicture}
\end{center}

\begin{apcnote}[AP scope: use it, do not derive it]
CED topic 8.9's exclusion: \emph{``Derivation of the
Henderson--Hasselbalch equation will not be assessed on the AP
Exam.''} So you will never be asked where it comes from --- only to
apply it. (For the curious: take $K_{a} =
[\ce{H+}][\ce{A-}]/[\ce{HA}]$, solve for $[\ce{H+}]$, and take
$-\log$ of both sides. That is the whole derivation, and it is not
examinable.)
\end{apcnote}

\begin{workedexample}[I do: three buffers of the same acid]
\textbf{Acetic acid has $K_{a} = \num{1.8e-5}$, so
p$K_{a} = 4.74$.}

\textbf{(a) Equal concentrations: \SI{0.10}{\Molar} \ce{CH3COOH} and
\SI{0.10}{\Molar} \ce{CH3COO-}.}
\[ \text{pH} = 4.74 + \log\!\left(\frac{0.10}{0.10}\right)
   = 4.74 + 0 = \mathbf{4.74} \]
When the two are equal, pH \emph{is} p$K_{a}$ --- the single most
useful fact in this chapter, and the basis of Ladder 12.

\textbf{(b) Twice as much base: \SI{0.10}{\Molar} \ce{HA},
\SI{0.20}{\Molar} \ce{A-}.}
\[ \text{pH} = 4.74 + \log 2 = 4.74 + 0.30 = \mathbf{5.04} \]

\textbf{(c) Using moles instead.} Suppose a solution contains
\SI{0.030}{\mole} \ce{HA} and \SI{0.020}{\mole} \ce{A-} in some volume:
\[ \text{pH} = 4.74 + \log\!\left(\frac{0.020}{0.030}\right)
   = 4.74 - 0.18 = \mathbf{4.56} \]
No volume was needed, because both species sit in the \emph{same}
solution and the volume divides out of the ratio. This is why the
titration ladders can work entirely in moles.

\textbf{Two errors to avoid.} The base goes on \textbf{top} --- putting
\ce{HA} there flips the sign of the correction and moves the pH the
wrong way. And the equation uses p$K_{a}$, not $K_{a}$: substituting
$\num{1.8e-5}$ where $4.74$ belongs produces nonsense you should catch
by inspection.

\textbf{When it stops being valid.} H--H assumes both species are
present in appreciable, comparable amounts. If the ratio is more
extreme than about 10:1 either way you no longer really have a buffer,
and the assumption that neither concentration changes much begins to
fail.
\end{workedexample}

\begin{yourturn}[YOUR TURN 3 --- four questions]
Acetic acid, p$K_{a} = 4.74$.
\begin{enumerate}[leftmargin=*,itemsep=0.5em]
  \item \SI{0.20}{\Molar} \ce{HA} and \SI{0.20}{\Molar} \ce{A-}. Find
        the pH. \workspace[1.6cm]
  \item \SI{0.10}{\Molar} \ce{HA} and \SI{0.010}{\Molar} \ce{A-}. Find
        the pH. \workspace[1.8cm]
  \item \SI{0.040}{\mole} \ce{A-} and \SI{0.020}{\mole} \ce{HA} in one
        solution. Find the pH. \workspace[1.8cm]
  \item Why may you use moles instead of concentrations?
        \workspace[1.6cm]
\end{enumerate}
\selfcheck{(a) 4.74 \quad (b) 3.74 \quad (c) 5.04 \quad (d) the volume
cancels in the ratio}
\keyonly{\begin{apcanswer}(a) Ratio 1, so pH $= \mathbf{4.74}$ ---
note the concentrations doubled from the worked example and the pH did
not move, because only the ratio matters. (b) $\log(0.010/0.10) =
\log(0.1) = -1$, so pH $= 4.74 - 1 = \mathbf{3.74}$. (c)
$\log(0.040/0.020) = \log 2 = 0.30$, so pH $= \mathbf{5.04}$. (d)
Because both species are dissolved in the \textbf{same solution and
therefore the same volume}, which divides out of the ratio
$[\ce{A-}]/[\ce{HA}]$.\end{apcanswer}}
\end{yourturn}

% =====================================================================
\section*{\sffamily Ladder 4 \textbullet\ Designing a buffer}
\zsec{15.2}

Given a target pH, choose the acid first and the ratio second.

\begin{workedexample}[I do: build a buffer at pH 5.00]
\textbf{Step 1 --- choose an acid whose p$K_{a}$ is close to the target
pH.} The correction term $\log(\text{ratio})$ should be small, which
means p$K_{a}$ within about one unit of the target. For pH 5.00, acetic
acid (p$K_{a} = 4.74$) is an excellent choice; \ce{HCN}
(p$K_{a} = 9.21$) would be a terrible one, because it would need a
ratio of about $10^{-4}$ and there would be almost no \ce{CN-} present
to do any buffering.

\textbf{Step 2 --- solve for the ratio.}
\[ 5.00 = 4.74 + \log\!\left(\frac{[\ce{A-}]}{[\ce{HA}]}\right) \]
\[ \log(\text{ratio}) = 0.26
   \quad\Longrightarrow\quad
   \text{ratio} = 10^{0.26} = \mathbf{1.8} \]

So you need about \textbf{1.8 moles of acetate for every 1 mole of
acetic acid}.

\textbf{Step 3 --- choose absolute amounts.} The ratio fixes the pH;
the \emph{concentrations} fix the capacity (Ladder 5). Using
\SI{0.10}{\Molar} \ce{HA} with \SI{0.18}{\Molar} \ce{A-} gives pH 5.00,
and so does \SI{1.0}{\Molar} with \SI{1.8}{\Molar} --- but the second
buffer will withstand ten times as much added acid or base.

\textbf{Two ways to make it in the lab.} Mix the weak acid with its
salt directly, or --- more common in practice --- take the weak acid and
add \emph{part} of the strong base needed to neutralise it. Half-neutralise
an acid and you get a buffer at exactly p$K_{a}$, which is Ladder 12.

\textbf{The choice-of-acid question is the assessed one.} If a question
gives you four acids and a target pH, the answer is the one whose
p$K_{a}$ is nearest --- and the justification is that a ratio near 1
gives comparable amounts of both species and therefore the greatest
capacity in both directions.
\end{workedexample}

\begin{yourturn}[YOUR TURN 4 --- four questions]
p$K_{a}$: acetic $4.74$, \ce{HF} $3.14$, \ce{HClO} $7.46$, \ce{HCN}
$9.21$.
\begin{enumerate}[leftmargin=*,itemsep=0.5em]
  \item Which acid would you choose for a pH 7.00 buffer? Why?
        \workspace[1.8cm]
  \item Which for a pH 3.00 buffer? \workspace[1.4cm]
  \item For an acid with p$K_{a} = 4.74$, what ratio gives pH 4.44?
        \workspace[1.8cm]
  \item Two buffers have the same ratio but one is ten times more
        concentrated. Do they have the same pH? \workspace[1.6cm]
\end{enumerate}
\selfcheck{(a) \ce{HClO} \quad (b) \ce{HF} \quad (c) 0.50 \quad (d)
yes}
\keyonly{\begin{apcanswer}(a) \textbf{\ce{HClO}}, p$K_{a} = 7.46$ ---
nearest to 7.00, so the ratio needed is close to 1 and both species are
present in comparable amounts. (b) \textbf{\ce{HF}}, p$K_{a} = 3.14$.
(c) $\log(\text{ratio}) = 4.44 - 4.74 = -0.30$, so ratio $=
10^{-0.30} = \mathbf{0.50}$ --- twice as much acid as base. (d)
\textbf{Yes.} pH depends only on the ratio. The concentrated one has a
much larger \emph{capacity} (Ladder 5), but the same
pH.\end{apcanswer}}
\end{yourturn}

% =====================================================================
\section*{\sffamily Ladder 5 \textbullet\ Buffer capacity}
\zsec{15.3}

Capacity is how much acid or base a buffer can absorb before it stops
working. It is set by the \emph{absolute} concentrations, not the
ratio.

\begin{center}
\begin{tikzpicture}[font=\sffamily\footnotesize]
  \node[font=\sffamily\small\bfseries] at (5.5,3.3)
    {two independent knobs};
  % second column at 4.0, not 3.4 -- "the concentrations" in bold sans
  % is wider than 2.8cm and ran into "set the capacity"
  \node[anchor=west,font=\sffamily\bfseries] at (0.6,2.65)
    {the \textbf{ratio}};
  \node[anchor=west] at (4.0,2.65)
    {sets the \textbf{pH} \quad(Ladder 3)};
  \node[anchor=west,font=\sffamily\bfseries] at (0.6,2.1)
    {the \textbf{concentrations}};
  \node[anchor=west] at (4.0,2.1)
    {set the \textbf{capacity}};
  \draw[apcgray] (0.5,1.75) -- (11.0,1.75);
  \node[anchor=west] at (0.6,1.3)
    {raise both concentrations, ratio fixed: \; same pH,
     \textbf{more} capacity};
  \node[anchor=west] at (0.6,0.75)
    {more conjugate \textbf{acid} than base: \; greater capacity for
     added \textbf{base}};
  \node[anchor=west] at (0.6,0.2)
    {more conjugate \textbf{base} than acid: \; greater capacity for
     added \textbf{acid}};
\end{tikzpicture}
\end{center}

\begin{workedexample}[I do: same pH, different staying power]
\textbf{Buffer A is \SI{0.10}{\Molar} in \ce{HA} and
\SI{0.10}{\Molar} in \ce{A-}. Buffer B is \SI{1.0}{\Molar} in each.
Compare them.}

\textbf{Same pH.} Both have ratio 1, so both sit at p$K_{a}$. The
Henderson--Hasselbalch equation cannot tell them apart.

\textbf{Very different capacity.} Buffer B contains ten times as many
moles of each component per litre, so it can absorb ten times as much
added acid or base before either component runs out. This is CED
8.10.A.1: raising the concentrations while holding the ratio constant
keeps the pH and increases the capacity.

\textbf{Now the asymmetry, which is 8.10.A.2 and less obvious.} A
buffer with \emph{more conjugate acid than base} has more \ce{HA}
available to react with added \ce{OH-} --- so it has greater capacity
for added \textbf{base}. A buffer with more conjugate base than acid has
more \ce{A-} to react with added \ce{H+}, so greater capacity for added
\textbf{acid}. The component in excess defends against the species it
reacts with.

\textbf{Which is worth saying slowly}, because the intuition often runs
backwards. Extra \emph{acid} in the buffer protects against added
\emph{base}, not against added acid --- because it is the acid
component that consumes hydroxide.

\textbf{And the practical consequence.} Capacity is greatest when the
two components are roughly equal \emph{and} both concentrated: equal
amounts give balanced protection in both directions, and high
concentration gives a lot of it. That is why buffers are normally made
at p$K_{a}$ and at the highest concentration the experiment tolerates.
\end{workedexample}

\begin{yourturn}[YOUR TURN 5 --- four questions]
\begin{enumerate}[leftmargin=*,itemsep=0.5em]
  \item Two buffers have ratio 1, one at \SI{0.05}{\Molar} and one at
        \SI{0.50}{\Molar}. Which has the higher pH?
        \workspace[1.6cm]
  \item Which has the greater capacity? \workspace[1.4cm]
  \item A buffer has more \ce{A-} than \ce{HA}. Is it better at
        absorbing added acid or added base? \workspace[1.8cm]
  \item What two conditions give the greatest overall capacity?
        \workspace[1.8cm]
\end{enumerate}
\selfcheck{(a) neither --- same pH \quad (b) the \SI{0.50}{\Molar} one
\quad (c) added acid \quad (d) equal amounts, both concentrated}
\keyonly{\begin{apcanswer}(a) \textbf{Neither --- they have the same
pH.} Ratio alone sets pH. (b) The \textbf{\SI{0.50}{\Molar}} buffer,
by a factor of ten: it contains ten times as many moles of each
component to consume added acid or base. (c) \textbf{Added acid.}
\ce{A-} is the component that reacts with \ce{H+}, and this buffer has
extra \ce{A-}. (d) \textbf{Roughly equal amounts of the two components}
(balanced protection in both directions) and \textbf{high
concentrations} of both (a lot of it).\end{apcanswer}}
\end{yourturn}

% =====================================================================
\section*{\sffamily Ladder 6 \textbullet\ Strong acid $+$ strong base}
\zsec{15.4}

The simplest mixture, and the template for all the others: do the
stoichiometry first, then look at what is left.

\begin{center}
\begin{tikzpicture}[font=\sffamily\footnotesize]
  \node[font=\sffamily\small\bfseries] at (5.5,2.9)
    {$\ce{H+ + OH- -> H2O}$ \;--- quantitative, one-way, complete};
  \foreach \y/\case/\what in {%
      2.2/{\ce{H+} in excess}/{pH from leftover $[\ce{H+}]$ and the
        \textbf{total} volume},
      1.65/{\ce{OH-} in excess}/{pOH from leftover $[\ce{OH-}]$, then
        pH $= 14 -$ pOH},
      1.1/{exactly equal}/{pH $= 7.00$ --- only water and spectator
        ions remain}} {
    \node[anchor=west,font=\sffamily\bfseries] at (0.6,\y) {\case};
    \node[anchor=west] at (3.6,\y) {\what};}
  \draw[apcgray] (0.5,2.55) -- (11.0,2.55);
  \draw[apcgray] (0.5,0.75) -- (11.0,0.75);
  \node[anchor=west] at (0.6,0.3)
    {the volume that matters is always the \textbf{combined} volume ---
     the commonest slip in this ladder};
\end{tikzpicture}
\end{center}

\begin{apcnote}[AP scope: a narrow exclusion, easy to over-read]
CED topic 8.9 excludes \emph{``computation of the change in pH
resulting from the addition of an acid or a base to a buffer''}. Read
that precisely: what is excluded is being handed an existing buffer and
asked how much the pH \emph{changes}. Computing the pH of a mixture
made by combining a weak acid with a strong base is topic
\textbf{8.4.A.2} and \textbf{is} assessed --- the CED says so
explicitly. Ladders 7--10 are that required material. The two
calculations look almost identical, so do not let the exclusion
frighten you off the titration work; it is the bulk of Unit 8's marks.
\end{apcnote}

\begin{workedexample}[I do: mix and find the leftover]
\textbf{Mix \SI{50.0}{\milli\litre} of \SI{0.100}{\Molar} \ce{HCl}
with \SI{30.0}{\milli\litre} of \SI{0.100}{\Molar} \ce{NaOH}. Find the
pH.}

\textbf{Step 1 --- moles, not concentrations.}
\[ n(\ce{H+}) = 0.0500 \times 0.100 = \SI{5.00e-3}{\mole} \]
\[ n(\ce{OH-}) = 0.0300 \times 0.100 = \SI{3.00e-3}{\mole} \]

\textbf{Step 2 --- react them completely.} They cancel one for one:
\[ \text{excess } \ce{H+} = 5.00 \times 10^{-3} - 3.00 \times 10^{-3}
   = \SI{2.00e-3}{\mole} \]

\textbf{Step 3 --- divide by the TOTAL volume.}
\[ V_{\text{total}} = 50.0 + 30.0 = \SI{80.0}{\milli\litre}
   = \SI{0.0800}{\litre} \]
\[ [\ce{H+}] = \frac{\num{2.00e-3}}{0.0800}
   = \SI{0.0250}{\Molar} \]

\textbf{Step 4 --- take the logarithm.}
\[ \text{pH} = -\log(0.0250) = \mathbf{1.60} \]

\textbf{The two habits this ladder is for.} Work in \textbf{moles}
through the stoichiometry, because concentrations change when you mix.
And divide by the \textbf{combined} volume at the end --- using
\SI{50.0}{\milli\litre} here would give pH 1.40 and be wrong by
0.2 units.

\textbf{Sanity check.} We started with excess acid, so the result must
be acidic; pH 1.60 is, and it is higher than the original
\SI{0.100}{\Molar} \ce{HCl} (pH 1.00) because some acid was neutralised
and the rest diluted. Both effects push the same way.
\end{workedexample}

\begin{yourturn}[YOUR TURN 6 --- four questions]
\begin{enumerate}[leftmargin=*,itemsep=0.5em]
  \item Mix \SI{25.0}{\milli\litre} of \SI{0.200}{\Molar} \ce{HCl}
        with \SI{25.0}{\milli\litre} of \SI{0.100}{\Molar} \ce{NaOH}.
        Find the excess in moles. \workspace[2.0cm]
  \item Find the pH of that mixture. \workspace[2.0cm]
  \item Mix equal volumes of \SI{0.100}{\Molar} \ce{HCl} and
        \SI{0.100}{\Molar} \ce{NaOH}. Give the pH.
        \workspace[1.6cm]
  \item Why must you use the combined volume?
        \workspace[1.6cm]
\end{enumerate}
\selfcheck{(a) \SI{2.50e-3}{\mole} \ce{H+} \quad (b) 1.30 \quad (c)
7.00 \quad (d) the solution is diluted on mixing}
\keyonly{\begin{apcanswer}(a) $n(\ce{H+}) = 0.0250 \times 0.200 =
\num{5.00e-3}$; $n(\ce{OH-}) = 0.0250 \times 0.100 = \num{2.50e-3}$;
excess $= \mathbf{\SI{2.50e-3}{\mole}}$ of \ce{H+}. (b) Total volume
\SI{50.0}{\milli\litre} $= \SI{0.0500}{\litre}$, so $[\ce{H+}] =
\num{2.50e-3}/0.0500 = \SI{0.0500}{\Molar}$ and pH $= \mathbf{1.30}$.
(c) Equal moles neutralise exactly, leaving only water and \ce{NaCl};
pH $= \mathbf{7.00}$. (d) Because \textbf{mixing dilutes both
solutions} --- the leftover moles are now spread through the sum of the
two volumes, and using either original volume gives a concentration
that is too high.\end{apcanswer}}
\end{yourturn}

% =====================================================================
\section*{\sffamily Ladder 7 \textbullet\ Weak acid $+$ strong base:
the stoichiometry step}
\zsec{15.4}

Now the reaction that generates a buffer. The strong base is consumed
\emph{completely}, so this step is stoichiometry --- no equilibrium
yet.

\[ \ce{HA + OH- -> A- + H2O} \]

\begin{center}
\begin{tikzpicture}[font=\sffamily\footnotesize]
  \node[font=\sffamily\small\bfseries] at (5.5,3.2)
    {the ICF table --- Initial, Change, Final, all in \textbf{moles}};
  \node[font=\sffamily\bfseries] at (2.9,2.6) {\ce{HA}};
  \node[font=\sffamily\bfseries] at (5.3,2.6) {\ce{OH-}};
  \node[font=\sffamily\bfseries] at (7.7,2.6) {\ce{A-}};
  \draw[apcgray] (0.9,2.3) -- (9.4,2.3);
  \foreach \y/\lab/\ha/\oh/\am in {%
      1.85/{\textbf{I}}/{$n_{\ce{HA}}$}/{$n_{\ce{OH-}}$}/{0},
      1.3/{\textbf{C}}/{$-n_{\ce{OH-}}$}/{$-n_{\ce{OH-}}$}/{$+n_{\ce{OH-}}$},
      0.75/{\textbf{F}}/{$n_{\ce{HA}} - n_{\ce{OH-}}$}/{\textbf{0}}/{$n_{\ce{OH-}}$}} {
    \node[anchor=west,font=\sffamily\bfseries] at (1.0,\y) {\lab};
    \node at (2.9,\y) {\ha};
    \node at (5.3,\y) {\oh};
    \node at (7.7,\y) {\am};}
  \draw[apcgray] (0.9,0.45) -- (9.4,0.45);
  \node[anchor=west] at (1.0,0.0)
    {\textbf{F, not E} --- the limiting reagent goes to zero, because a
     strong base reacts completely};
\end{tikzpicture}
\end{center}

\begin{apctrap}
\textbf{It is an IC\emph{F} table, not an IC\emph{E} table.} The strong
base is not in equilibrium with anything --- it is consumed to
\textbf{zero}. Writing ``$-x$'' here, or leaving leftover \ce{OH-} in an
equilibrium expression, confuses two different steps. Do the
stoichiometry to completion first; \emph{then}, if a weak species
remains, do the equilibrium.
\end{apctrap}

\begin{workedexample}[I do: what is in the flask?]
\textbf{\SI{50.0}{\milli\litre} of \SI{0.100}{\Molar} acetic acid is
treated with \SI{20.0}{\milli\litre} of \SI{0.100}{\Molar} \ce{NaOH}.
What is present afterwards?}

\textbf{Moles first.}
\[ n(\ce{HA}) = 0.0500 \times 0.100 = \SI{5.00e-3}{\mole} \]
\[ n(\ce{OH-}) = 0.0200 \times 0.100 = \SI{2.00e-3}{\mole} \]

\textbf{The ICF table.}
\begin{center}
\renewcommand{\arraystretch}{1.15}
\begin{tabular}{@{}lccc@{}}
\toprule
(mol) & \ce{HA} & \ce{OH-} & \ce{A-} \\
\midrule
\textbf{I} & \num{5.00e-3} & \num{2.00e-3} & 0 \\
\textbf{C} & $-\num{2.00e-3}$ & $-\num{2.00e-3}$ & $+\num{2.00e-3}$ \\
\textbf{F} & \num{3.00e-3} & \textbf{0} & \num{2.00e-3} \\
\bottomrule
\end{tabular}
\end{center}

\textbf{Read the F row and name the situation.} Both \ce{HA} and
\ce{A-} are present in appreciable amounts --- \textbf{this is a
buffer}. The hydroxide is gone.

\textbf{So the pH comes from Henderson--Hasselbalch}, using the mole
amounts directly (Ladder 3 showed the volume cancels):
\[ \text{pH} = 4.74 + \log\!\left(\frac{\num{2.00e-3}}{\num{3.00e-3}}\right)
   = 4.74 - 0.18 = \mathbf{4.56} \]

\textbf{The F row tells you which method to use}, and that is the real
skill:
\begin{itemize}[leftmargin=1.4em,itemsep=0.15em]
  \item \ce{HA} and \ce{A-} both present $\Rightarrow$ buffer, use H--H
  \item only \ce{A-} left $\Rightarrow$ equivalence point, use $K_{b}$
        (Ladder 9)
  \item \ce{OH-} left over $\Rightarrow$ past equivalence, use excess
        hydroxide (Ladder 8)
\end{itemize}
Every weak-acid titration question is one of these three, and the ICF
table decides which.
\end{workedexample}

\begin{yourturn}[YOUR TURN 7 --- four questions]
\SI{50.0}{\milli\litre} of \SI{0.100}{\Molar} \ce{HA}
($\text{p}K_{a} = 4.74$) is treated with \SI{0.100}{\Molar} \ce{NaOH}.
\begin{enumerate}[leftmargin=*,itemsep=0.5em]
  \item After \SI{10.0}{\milli\litre} of base, give the moles of
        \ce{HA} and \ce{A-}. \workspace[2.0cm]
  \item Find the pH at that point. \workspace[1.8cm]
  \item After \SI{25.0}{\milli\litre}, give the moles of each and the
        pH. \workspace[2.0cm]
  \item Why is it an ICF and not an ICE table?
        \workspace[1.6cm]
\end{enumerate}
\selfcheck{(a) \SI{4.00e-3}{} and \SI{1.00e-3}{\mole} \quad (b) 4.14
\quad (c) \SI{2.50e-3}{} each, pH 4.74 \quad (d) the strong base is
consumed completely}
\keyonly{\begin{apcanswer}(a) $n(\ce{OH-}) = \num{1.00e-3}$, so
\ce{HA} $= 5.00 - 1.00 = \mathbf{\SI{4.00e-3}{\mole}}$ and \ce{A-}
$= \mathbf{\SI{1.00e-3}{\mole}}$. (b) pH $= 4.74 +
\log(1.00/4.00) = 4.74 - 0.60 = \mathbf{4.14}$. (c) $n(\ce{OH-}) =
\num{2.50e-3}$, leaving \ce{HA} $= \num{2.50e-3}$ and \ce{A-} $=
\num{2.50e-3}$ --- equal, so pH $= \mathbf{4.74} = \text{p}K_{a}$
(Ladder 12). (d) Because the \textbf{strong base reacts completely} ---
it is consumed to zero rather than reaching an equilibrium. The
equilibrium step comes afterwards, on whatever weak species
remain.\end{apcanswer}}
\end{yourturn}

% =====================================================================
\section*{\sffamily Ladder 8 \textbullet\ The three regions of a
titration}
\zsec{15.4}

One titration, three different calculations, and the ICF table tells
you which you are in.

\begin{center}
\begin{tikzpicture}
\begin{axis}[
  width=11.4cm, height=6.6cm,
  xlabel={volume of \SI{0.100}{\Molar} \ce{NaOH} added
    (\si{\milli\litre})},
  ylabel={pH},
  xmin=0, xmax=70, ymin=0, ymax=14,
  xtick={0,10,25,50,60,70}, ytick={0,2,4,6,8,10,12,14},
  grid=major, grid style={apcgray!25},
  label style={font=\footnotesize}, tick label style={font=\footnotesize},
]
\addplot[seriesA, very thick] coordinates {
  (0,2.87) (10,4.14) (20,4.56) (25,4.74) (30,4.92) (40,5.34)
  (45,5.69) (49,6.43) (50,8.72) (51,11.00) (55,11.68) (60,11.96)
  (70,12.22)};
\addplot[only marks, mark=*, mark size=1.8pt, black]
  coordinates {(25,4.74) (50,8.72)};
\draw[densely dashed, apcgray] (axis cs:50,0) -- (axis cs:50,8.72);
\draw[densely dashed, apcgray] (axis cs:25,0) -- (axis cs:25,4.74);
\node[font=\footnotesize,anchor=west] at (axis cs:2,7.2)
  {\textbf{buffer region}};
\node[font=\footnotesize,anchor=west] at (axis cs:2,6.4)
  {use H--H};
\node[font=\footnotesize,anchor=west] at (axis cs:52,6.0)
  {equivalence};
\node[font=\footnotesize,anchor=west] at (axis cs:52,5.2)
  {pH 8.72};
\node[font=\footnotesize,anchor=west] at (axis cs:27,3.0)
  {half-equivalence, pH $=$ p$K_{a}$};
\node[font=\footnotesize,anchor=east] at (axis cs:68,10.2)
  {excess \ce{OH-}};
\end{axis}
\end{tikzpicture}
\end{center}

\begin{workedexample}[I do: past the equivalence point]
\textbf{Continue the titration from Ladder 7 to
\SI{60.0}{\milli\litre} of \SI{0.100}{\Molar} \ce{NaOH}. Find the pH.}

\textbf{ICF first, as always.}
\[ n(\ce{HA}) = \num{5.00e-3}
   \qquad
   n(\ce{OH-}) = 0.0600 \times 0.100 = \num{6.00e-3} \]
The base is now in excess:
\[ \text{excess } \ce{OH-} = 6.00 \times 10^{-3} - 5.00 \times 10^{-3}
   = \SI{1.00e-3}{\mole} \]

\textbf{Read the F row.} There is leftover \emph{strong} base. Strong
base overwhelms the weak conjugate base \ce{A-} that is also present ---
by several orders of magnitude --- so the acetate contributes nothing
worth counting and the pH comes from the excess hydroxide alone.

\textbf{Divide by the total volume and finish.}
\[ V = 50.0 + 60.0 = \SI{110.0}{\milli\litre} = \SI{0.110}{\litre} \]
\[ [\ce{OH-}] = \frac{\num{1.00e-3}}{0.110}
   = \SI{9.09e-3}{\Molar} \]
\[ \text{pOH} = 2.04
   \qquad
   \text{pH} = 14.00 - 2.04 = \mathbf{11.96} \]

\textbf{Notice this is Ladder 6's calculation.} Once strong base is in
excess, a weak-acid titration behaves exactly like a strong-acid one ---
which is why the two curves in Ladder 11 converge after the equivalence
point.

\textbf{The three regions, one last time.} Before equivalence: buffer,
use H--H. At equivalence: only \ce{A-}, use its $K_{b}$ (Ladder 9).
After equivalence: excess \ce{OH-}, use it directly. Identify the region
from the ICF table and the method chooses itself.
\end{workedexample}

\begin{yourturn}[YOUR TURN 8 --- four questions]
\SI{50.0}{\milli\litre} of \SI{0.100}{\Molar} \ce{HA}
(p$K_{a} = 4.74$) with \SI{0.100}{\Molar} \ce{NaOH}.
\begin{enumerate}[leftmargin=*,itemsep=0.5em]
  \item Which region is \SI{40.0}{\milli\litre} in, and what is the pH?
        \workspace[2.0cm]
  \item Which region is \SI{55.0}{\milli\litre} in?
        \workspace[1.4cm]
  \item Find the pH at \SI{55.0}{\milli\litre}.
        \workspace[2.2cm]
  \item Why does the acetate present after equivalence not affect the
        pH? \workspace[1.8cm]
\end{enumerate}
\selfcheck{(a) buffer, pH 5.34 \quad (b) past equivalence \quad (c)
11.68 \quad (d) the strong base overwhelms it}
\keyonly{\begin{apcanswer}(a) \textbf{Buffer region}: \ce{HA} $=
\num{1.00e-3}$, \ce{A-} $= \num{4.00e-3}$, so pH $= 4.74 + \log 4 =
4.74 + 0.60 = \mathbf{5.34}$. (b) \textbf{Past the equivalence point}
--- \SI{50.0}{\milli\litre} was needed to neutralise the acid. (c)
Excess \ce{OH-} $= \num{5.50e-3} - \num{5.00e-3} = \num{5.00e-4}$ mol
in \SI{0.105}{\litre}, so $[\ce{OH-}] = \num{4.76e-3}$, pOH $= 2.32$,
pH $= \mathbf{11.68}$. (d) Because \ce{A-} is a \textbf{very weak
base} ($K_{b} = \num{5.6e-10}$) and the leftover \ce{OH-} is a
\emph{strong} base present at a concentration millions of times larger
in effect. The weak contribution is negligible beside
it.\end{apcanswer}}
\end{yourturn}

% =====================================================================
\section*{\sffamily Ladder 9 \textbullet\ pH at the equivalence point}
\zsec{15.4}

At equivalence the acid is exactly consumed and \emph{only the
conjugate base remains}. That is a weak-base problem --- Chapter 14,
Ladder 11.

\begin{apctrap}
\textbf{The equivalence point of a weak acid titration is not pH 7.} It
is \emph{basic}, because the solution at that moment is a solution of
\ce{A-}, which is a weak base. Answering ``7'' is the standard error and
costs the whole part. Only a \emph{strong} acid with a strong base
gives 7.
\end{apctrap}

\begin{workedexample}[I do: the equivalence point, in full]
\textbf{Find the pH at the equivalence point of the Ladder 7 titration:
\SI{50.0}{\milli\litre} of \SI{0.100}{\Molar} acetic acid with
\SI{0.100}{\Molar} \ce{NaOH}. $K_{a} = \num{1.8e-5}$.}

\textbf{Step 1 --- how much base, and what is left?} Equivalence needs
\SI{5.00e-3}{\mole} of \ce{OH-}, which is
\SI{50.0}{\milli\litre}. The F row is: no \ce{HA}, no \ce{OH-}, and
\SI{5.00e-3}{\mole} of \ce{A-}.

\textbf{Step 2 --- concentration of the conjugate base}, using the
\textbf{total} volume:
\[ [\ce{A-}] = \frac{\num{5.00e-3}}{0.100\;\si{\litre}}
   = \SI{0.0500}{\Molar} \]
The dilution is real --- the acetate sits in twice the original volume.

\textbf{Step 3 --- treat it as the weak base it is.}
\[ K_{b} = \frac{K_{w}}{K_{a}}
   = \frac{\num{1.0e-14}}{\num{1.8e-5}} = \num{5.6e-10} \]
\[ x = \sqrt{K_{b}c} = \sqrt{(\num{5.6e-10})(0.0500)}
     = \SI{5.3e-6}{\Molar} = [\ce{OH-}] \]
Check: $\num{5.3e-6}/0.0500$ is far below 5\%. \checkmark

\textbf{Step 4 --- convert to pH}, remembering $x$ is $[\ce{OH-}]$:
\[ \text{pOH} = 5.28
   \qquad
   \text{pH} = 14.00 - 5.28 = \mathbf{8.72} \]

\textbf{Basic, as it must be.} And note the two dilution traps this
problem contains: the volume doubled, and $x$ gave pOH rather than pH.
Either slip alone produces a wrong answer that still \emph{looks}
plausible.

\textbf{The four-step shape is worth memorising} --- moles to
equivalence, concentration in the total volume, $K_{b}$ from $K_{a}$,
then the weak-base ICE. It is the same shape for a weak base titrated
with strong acid (Ladder 10), with acid and base interchanged.
\end{workedexample}

\begin{yourturn}[YOUR TURN 9 --- four questions]
\begin{enumerate}[leftmargin=*,itemsep=0.5em]
  \item Is the equivalence pH of a weak acid $+$ strong base titration
        above or below 7? Why? \workspace[1.8cm]
  \item What volume of \SI{0.100}{\Molar} \ce{NaOH} reaches
        equivalence for \SI{25.0}{\milli\litre} of \SI{0.100}{\Molar}
        \ce{HA}? \workspace[1.8cm]
  \item At that equivalence point, give $[\ce{A-}]$.
        \workspace[1.8cm]
  \item Which constant do you use at the equivalence point, $K_{a}$ or
        $K_{b}$? \workspace[1.6cm]
\end{enumerate}
\selfcheck{(a) above 7 --- \ce{A-} is a weak base \quad (b)
\SI{25.0}{\milli\litre} \quad (c) \SI{0.0500}{\Molar} \quad (d)
$K_{b}$}
\keyonly{\begin{apcanswer}(a) \textbf{Above 7.} At equivalence the only
solute is \ce{A-}, the conjugate base of a weak acid, and it takes a
proton from water to give \ce{OH-}. (b)
\textbf{\SI{25.0}{\milli\litre}} --- equal concentrations means equal
volumes. (c) $n(\ce{A-}) = \num{2.50e-3}$ mol in a total
\SI{50.0}{\milli\litre}, so $[\ce{A-}] =
\mathbf{\SI{0.0500}{\Molar}}$. (d) \textbf{$K_{b}$}, obtained as
$K_{w}/K_{a}$ --- the species present is a base, so you need a base
constant.\end{apcanswer}}
\end{yourturn}

% =====================================================================
\section*{\sffamily Ladder 10 \textbullet\ Weak base $+$ strong acid}
\zsec{15.4}

The mirror image. Same three regions, same two-step method, everything
reflected about pH 7.

\[ \ce{B + H+ -> BH+} \]

\begin{workedexample}[I do: ammonia titrated with hydrochloric acid]
\textbf{\SI{50.0}{\milli\litre} of \SI{0.100}{\Molar} \ce{NH3}
($K_{b} = \num{1.8e-5}$) is titrated with \SI{0.100}{\Molar} \ce{HCl}.}

\textbf{First convert to the acid constant you will need.} The buffer
here is \ce{NH3}/\ce{NH4+}, and Henderson--Hasselbalch is written for a
conjugate acid:
\[ K_{a}(\ce{NH4+}) = \frac{\num{1.0e-14}}{\num{1.8e-5}}
   = \num{5.6e-10}
   \qquad
   \text{p}K_{a} = 9.26 \]

\textbf{At half-equivalence (\SI{25.0}{\milli\litre} of acid):} the ICF
row gives \SI{2.50e-3}{\mole} each of \ce{NH3} and \ce{NH4+}, so
\[ \text{pH} = 9.26 + \log 1 = \mathbf{9.26} \]
Still basic --- the buffer region of a weak-base titration sits above 7.

\textbf{At equivalence (\SI{50.0}{\milli\litre}):} everything is
\ce{NH4+}, at
\[ [\ce{NH4+}] = \frac{\num{5.00e-3}}{0.100} = \SI{0.0500}{\Molar} \]
\ce{NH4+} is a weak \emph{acid}, so
\[ x = \sqrt{(\num{5.6e-10})(0.0500)} = \num{5.3e-6} = [\ce{H+}] \]
\[ \text{pH} = \mathbf{5.28} \]

\textbf{Acidic at equivalence}, which is the mirror of Ladder 9's 8.72
--- and notice $8.72 + 5.28 = 14.00$ exactly, because acetic acid's
$K_{a}$ and ammonia's $K_{b}$ happen to be the same number. That
symmetry is a good check when a problem is built from that pair.

\textbf{The rule for equivalence-point pH, in one line.} Strong$+$strong
gives 7; weak acid $+$ strong base gives above 7; weak base $+$ strong
acid gives below 7. The \emph{weak} partner determines which way, and
its conjugate is what remains in the flask.
\end{workedexample}

\begin{yourturn}[YOUR TURN 10 --- four questions]
\ce{NH3}, $K_{b} = \num{1.8e-5}$, p$K_{a}(\ce{NH4+}) = 9.26$.
\begin{enumerate}[leftmargin=*,itemsep=0.5em]
  \item Is the equivalence pH above or below 7? Why?
        \workspace[1.8cm]
  \item In the buffer region, which two species are present?
        \workspace[1.6cm]
  \item \SI{50.0}{\milli\litre} of \SI{0.100}{\Molar} \ce{NH3} plus
        \SI{20.0}{\milli\litre} of \SI{0.100}{\Molar} \ce{HCl}: find
        the pH. \workspace[2.2cm]
  \item Which constant do you need for the H--H equation here, and why?
        \workspace[1.8cm]
\end{enumerate}
\selfcheck{(a) below 7 --- \ce{NH4+} is a weak acid \quad (b) \ce{NH3}
and \ce{NH4+} \quad (c) 9.44 \quad (d) $K_{a}$ of \ce{NH4+}}
\keyonly{\begin{apcanswer}(a) \textbf{Below 7}: at equivalence the only
solute is \ce{NH4+}, the conjugate acid of a weak base, which donates a
proton to water. (b) \textbf{\ce{NH3} and \ce{NH4+}} --- a conjugate
pair. (c) $n(\ce{NH3}) = \num{5.00e-3}$, $n(\ce{H+}) = \num{2.00e-3}$,
so \ce{NH3} $= \num{3.00e-3}$ and \ce{NH4+} $= \num{2.00e-3}$;
pH $= 9.26 + \log(3.00/2.00) = 9.26 + 0.18 = \mathbf{9.44}$. Note the
\emph{base} goes on top, and here the base is \ce{NH3}. (d)
\textbf{$K_{a}$ of \ce{NH4+}}, obtained as $K_{w}/K_{b}$ ---
Henderson--Hasselbalch is written in terms of the conjugate
\emph{acid}'s p$K_{a}$.\end{apcanswer}}
\end{yourturn}

% =====================================================================
\section*{\sffamily Ladder 11 \textbullet\ Reading titration curves}
\zsec{15.4--15.5}

Two curves, one titrant, and the differences are all diagnostic.

\begin{center}
\begin{tikzpicture}
\begin{axis}[
  width=11.4cm, height=6.4cm,
  xlabel={volume of \SI{0.100}{\Molar} \ce{NaOH} (\si{\milli\litre})},
  ylabel={pH},
  xmin=0, xmax=70, ymin=0, ymax=14,
  xtick={0,10,20,30,40,50,60,70}, ytick={0,2,4,6,8,10,12,14},
  grid=major, grid style={apcgray!25},
  legend pos=north west, legend style={font=\footnotesize,
    cells={anchor=west}},
  label style={font=\footnotesize}, tick label style={font=\footnotesize},
]
\addplot[seriesA, very thick] coordinates {
  (0,2.87) (10,4.14) (20,4.56) (25,4.74) (30,4.92) (40,5.34)
  (45,5.69) (49,6.43) (50,8.72) (51,11.00) (55,11.68) (60,11.96)
  (70,12.22)};
\addlegendentry{weak acid (\ce{CH3COOH})}
\addplot[seriesB, very thick] coordinates {
  (0,1.00) (10,1.18) (25,1.48) (40,1.95) (49,3.00) (50,7.00)
  (51,11.00) (60,11.96) (70,12.22)};
\addlegendentry{strong acid (\ce{HCl})}
\addplot[only marks, mark=*, mark size=1.6pt, black]
  coordinates {(50,8.72) (50,7.00)};
\end{axis}
\end{tikzpicture}
\end{center}

\begin{workedexample}[I do: four things a curve tells you]
\textbf{Compare the two curves above. Both titrate
\SI{50.0}{\milli\litre} of a \SI{0.100}{\Molar} acid with the same
base.}

\textbf{1. The starting pH.} The weak acid starts at 2.87, the strong
acid at 1.00. A high starting pH for a given concentration is the first
sign of a weak acid.

\textbf{2. The buffer plateau.} The weak-acid curve has a long, gentle
stretch between about 10 and \SI{40}{\milli\litre} where the pH barely
moves. The strong-acid curve has no such region --- there is no
conjugate pair to buffer it. \textbf{A flat middle region means a weak
acid.}

\textbf{3. The equivalence pH.} 8.72 for the weak acid, 7.00 for the
strong. Above 7 identifies a weak acid immediately.

\textbf{4. The equivalence \emph{volume} is the same for both:}
\SI{50.0}{\milli\litre}. This is the point students find surprising ---
the volume depends only on \emph{moles} of acid, not on strength. A weak
acid needs exactly as much base as a strong one at the same
concentration. Strength changes the \emph{shape} of the curve, not
where the equivalence point falls.

\textbf{And after equivalence the curves coincide}, because both
solutions are then just excess \ce{NaOH} diluted into a similar volume.
Whatever acid you started with stops mattering once it is gone.

\textbf{The steep jump is what makes titration work as a measurement.}
Near equivalence, a fraction of a millilitre swings the pH by several
units, so the endpoint is sharp and easy to detect --- which is Ladder
13's business.
\end{workedexample}

\begin{yourturn}[YOUR TURN 11 --- four questions]
\begin{enumerate}[leftmargin=*,itemsep=0.5em]
  \item Two acids of the same concentration are titrated. One curve
        starts at pH 1.0 and one at 3.0. Which is weak?
        \workspace[1.6cm]
  \item A titration curve has equivalence pH 9.0. What kind of acid was
        it? \workspace[1.6cm]
  \item Does a weak acid need more, less, or the same volume of base to
        reach equivalence as a strong acid of the same concentration?
        \workspace[1.8cm]
  \item What feature of the curve shows a buffer region?
        \workspace[1.6cm]
\end{enumerate}
\selfcheck{(a) the one starting at 3.0 \quad (b) weak \quad (c) the
same \quad (d) the flat, gently sloping stretch}
\keyonly{\begin{apcanswer}(a) The one starting at \textbf{pH 3.0} --- a
weak acid ionises only partly, so $[\ce{H+}]$ is much lower at the same
nominal concentration. (b) A \textbf{weak} acid: equivalence above 7
means the conjugate base left in the flask is appreciably basic. (c)
\textbf{The same volume.} Equivalence depends only on the \emph{moles}
of acid present, not on how completely it ionises. (d) The
\textbf{flat, gently sloping stretch} before the equivalence point,
where added base changes the pH very little.\end{apcanswer}}
\end{yourturn}

% =====================================================================
\section*{\sffamily Ladder 12 \textbullet\ The half-equivalence point}
\zsec{15.4}

The single most useful point on any weak-acid curve.

\begin{center}
\begin{tikzpicture}[font=\sffamily\footnotesize]
  \node[font=\sffamily\small\bfseries] at (5.5,2.8)
    {at half-equivalence, exactly half the acid has been converted};
  \node[anchor=west] at (0.6,2.15)
    {$n(\ce{HA}) = n(\ce{A-})$ \quad$\Longrightarrow$\quad ratio $=1$
     \quad$\Longrightarrow$\quad $\log 1 = 0$};
  \draw[-{Stealth[length=2.5mm]},thick] (5.5,1.85) -- (5.5,1.45);
  \node[font=\sffamily\bfseries] at (5.5,1.1)
    {pH $=$ p$K_{a}$};
  \draw[apcgray] (0.5,0.7) -- (11.0,0.7);
  \node[anchor=west] at (0.6,0.25)
    {so you can \textbf{read $K_{a}$ off a titration curve} --- which is
     how weak-acid constants are measured};
\end{tikzpicture}
\end{center}

\begin{workedexample}[I do: measuring $K_{a}$ from a graph]
\textbf{A weak acid is titrated. The equivalence point is at
\SI{40.0}{\milli\litre} of base, and the pH at
\SI{20.0}{\milli\litre} is 5.20. Find $K_{a}$.}

\textbf{Locate half-equivalence.} Half of \SI{40.0}{\milli\litre} is
\SI{20.0}{\milli\litre} --- so the given point \emph{is} the
half-equivalence point.

\textbf{Apply the identity.} There, $[\ce{HA}] = [\ce{A-}]$, so the
logarithm vanishes and
\[ \text{pH} = \text{p}K_{a} = 5.20 \]

\textbf{Convert.}
\[ K_{a} = 10^{-5.20} = \mathbf{\num{6.3e-6}} \]

\textbf{This is a real laboratory method}, not a textbook trick. Titrate
a weak acid, find the equivalence volume from the steep jump, halve it,
read the pH there, and you have p$K_{a}$ --- no separate experiment
required.

\textbf{Why the point is flat, which is the practical payoff.} At
half-equivalence the curve is at its most horizontal, so a small error
in reading the volume produces almost no error in the pH. The
measurement is unusually forgiving, which is exactly what you want from
a method.

\textbf{Two cautions.} Half-equivalence is half the volume \emph{to
equivalence}, not half the total volume added, and not the middle of the
graph. And this identity holds only for a weak acid or base --- a strong
acid has no buffer region and no p$K_{a}$ to read.
\end{workedexample}

\begin{yourturn}[YOUR TURN 12 --- four questions]
\begin{enumerate}[leftmargin=*,itemsep=0.5em]
  \item Equivalence is at \SI{30.0}{\milli\litre}. Where is
        half-equivalence? \workspace[1.4cm]
  \item The pH there is 4.20. Give p$K_{a}$ and $K_{a}$.
        \workspace[1.8cm]
  \item Why does the logarithm term vanish at half-equivalence?
        \workspace[1.8cm]
  \item Can you use this method on a strong acid? Explain.
        \workspace[1.8cm]
\end{enumerate}
\selfcheck{(a) \SI{15.0}{\milli\litre} \quad (b) 4.20 and
\num{6.3e-5} \quad (c) the ratio is 1 \quad (d) no --- no buffer
region}
\keyonly{\begin{apcanswer}(a) \textbf{\SI{15.0}{\milli\litre}}.
(b) p$K_{a} = \mathbf{4.20}$, so $K_{a} = 10^{-4.20} =
\mathbf{\num{6.3e-5}}$. (c) Because exactly half the acid has been
converted, so $[\ce{HA}] = [\ce{A-}]$, the ratio is 1, and
$\log 1 = 0$. (d) \textbf{No.} A strong acid ionises completely, so
there is no \ce{HA}/\ce{A-} conjugate pair, no buffer region, and no
$K_{a}$ to measure --- its curve has no flat stretch at
all.\end{apcanswer}}
\end{yourturn}

% =====================================================================
\section*{\sffamily Ladder 13 \textbullet\ Choosing an indicator}
\zsec{15.5}

An indicator is itself a weak acid whose two forms differ in colour. It
changes over roughly $\text{p}K_{a} \pm 1$.

\begin{center}
\begin{tikzpicture}[font=\sffamily\footnotesize]
  \draw[-{Stealth[length=3mm]},very thick] (0.7,0.6) -- (10.9,0.6);
  \foreach \x/\v in {1.0/0, 2.4/2, 3.8/4, 5.2/6, 6.6/8, 8.0/10,
                     9.4/12, 10.5/14} {
    \fill (\x,0.6) circle (0.06);
    \node[font=\sffamily\footnotesize] at (\x,0.25) {\v};}
  \node[anchor=north,font=\sffamily\footnotesize] at (5.8,-0.15) {pH};
  % Names sit BESIDE their bands, not inside them: each name is wider
  % than the band it labels (a 2 pH-unit range is only ~1.4cm here), so
  % centring the text on the band spilled it over both edges.
  \fill[apcgray!45] (3.24,1.0) rectangle (4.71,1.45);
  \node[anchor=west,font=\sffamily\footnotesize] at (4.85,1.22)
    {methyl red \;\textbullet\; 4.2--6.3};
  \fill[apcgray!45] (4.46,1.65) rectangle (5.58,2.1);
  \node[anchor=west,font=\sffamily\footnotesize] at (5.72,1.87)
    {bromothymol blue \;\textbullet\; 6.0--7.6};
  \fill[apcgray!45] (6.0,2.3) rectangle (7.7,2.75);
  \node[anchor=west,font=\sffamily\footnotesize] at (7.84,2.52)
    {phenolphthalein \;\textbullet\; 8.3--10.0};
  \draw[apcgray] (0.6,3.05) -- (11.0,3.05);
  \node[anchor=west,font=\sffamily\bfseries] at (0.7,3.45)
    {pick the indicator whose range brackets the \emph{equivalence} pH};
\end{tikzpicture}
\end{center}

\begin{workedexample}[I do: three titrations, three indicators]
\textbf{The rule: the indicator's colour-change range must contain the
equivalence pH, so the colour change coincides with the steep jump.}

\textbf{Strong acid $+$ strong base}, equivalence pH 7.00.
\textbf{Bromothymol blue} (6.0--7.6) brackets it neatly. Phenolphthalein
also works in practice, because the jump at equivalence is so nearly
vertical that it passes through 8.3 within a fraction of a drop.

\textbf{Weak acid $+$ strong base}, equivalence pH 8.72 (Ladder 9).
\textbf{Phenolphthalein} (8.3--10.0) is the correct choice. Methyl red
would change colour around pH 5 --- long before equivalence --- and
would report a badly low volume.

\textbf{Weak base $+$ strong acid}, equivalence pH 5.28 (Ladder 10).
\textbf{Methyl red} (4.2--6.3). Phenolphthalein would change far too
late.

\textbf{Why the steep jump makes this work at all.} Near equivalence the
pH moves several units within a fraction of a millilitre, so any
indicator whose range lies inside that jump changes colour at
essentially the right volume. The steeper the jump, the more forgiving
the choice --- which is why a strong--strong titration tolerates almost
any indicator and a weak--weak titration is not worth attempting.

\textbf{Endpoint versus equivalence point.} The \textbf{equivalence
point} is where the moles match --- a fact about the chemistry. The
\textbf{endpoint} is where the indicator changes --- what you actually
observe. A well-chosen indicator makes them nearly coincide; a poor one
makes them differ, and that difference is systematic experimental error.
\end{workedexample}

\begin{yourturn}[YOUR TURN 13 --- four questions]
\begin{enumerate}[leftmargin=*,itemsep=0.5em]
  \item Which indicator suits an equivalence pH of 9.0?
        \workspace[1.6cm]
  \item Which suits an equivalence pH of 5.0?
        \workspace[1.6cm]
  \item Distinguish the endpoint from the equivalence point.
        \workspace[1.8cm]
  \item Why can a strong--strong titration tolerate almost any
        indicator? \workspace[1.8cm]
\end{enumerate}
\selfcheck{(a) phenolphthalein \quad (b) methyl red \quad (c) observed
colour change vs where moles match \quad (d) the pH jump is very steep}
\keyonly{\begin{apcanswer}(a) \textbf{Phenolphthalein} (8.3--10.0).
(b) \textbf{Methyl red} (4.2--6.3). (c) The \textbf{equivalence point}
is where the moles of acid and base are stoichiometrically equal --- a
property of the chemistry. The \textbf{endpoint} is where the indicator
visibly changes colour --- what the experimenter observes. A good
indicator makes them nearly coincide. (d) Because the pH \textbf{jumps
almost vertically} through several units within a fraction of a drop, so
any indicator changing anywhere in that range does so at essentially the
correct volume.\end{apcanswer}}
\end{yourturn}

% =====================================================================
\section*{\sffamily Ladder 14 \textbullet\ Polyprotic titration curves}
\zsec{15.6}

A diprotic acid gives up its protons in stages, so its curve has
\textbf{two} equivalence points.

\begin{center}
\begin{tikzpicture}
\begin{axis}[
  width=11.0cm, height=5.8cm,
  xlabel={volume of strong base added}, ylabel={pH},
  xmin=0, xmax=100, ymin=0, ymax=14,
  xtick={0,25,50,75,100},
  xticklabels={0,{$\frac{1}{2}$eq},{1st eq},{$\frac{3}{2}$eq},
    {2nd eq}},
  ytick={0,2,4,6,8,10,12,14},
  grid=major, grid style={apcgray!25},
  label style={font=\footnotesize}, tick label style={font=\footnotesize},
]
\addplot[seriesA, very thick] coordinates {
  (0,1.9) (10,2.6) (25,3.1) (40,3.7) (48,4.6) (50,5.6) (52,6.6)
  (60,7.0) (75,7.2) (88,8.4) (95,9.6) (98,11.0) (100,11.6)};
\addplot[only marks, mark=*, mark size=1.6pt, black]
  coordinates {(25,3.1) (50,5.6) (75,7.2) (100,11.6)};
\node[font=\footnotesize,anchor=west] at (axis cs:3,10.6)
  {two buffer regions, two jumps};
\node[font=\footnotesize,anchor=west] at (axis cs:3,9.6)
  {pH $=$ p$K_{a1}$ and p$K_{a2}$ at the half-points};
\end{axis}
\end{tikzpicture}
\end{center}

\begin{apcnote}[AP scope: shapes yes, species computations no]
CED topic 8.5's exclusion: \emph{``Computation of the concentration of
each species present in the titration curve for polyprotic acids will
not be assessed on the AP Exam.''} The same sentence goes on to say
that these computations \emph{are} in scope for \textbf{monoprotic}
acids (that is Ladders 7--10), \emph{``as is qualitative reasoning
regarding what species are present''}. So for polyprotic acids: read
the curve, identify the equivalence points, say which species dominates
where --- but do not expect to calculate concentrations.
\end{apcnote}

\begin{workedexample}[I do: reading a diprotic curve]
\textbf{A diprotic acid \ce{H2A} is titrated with strong base. What
happens, in order?}

\textbf{First stage.} $\ce{H2A + OH- -> HA- + H2O}$ consumes the first
proton. The buffer region here is \ce{H2A}/\ce{HA-}, and at its
half-point pH $=$ p$K_{a1}$.

\textbf{First equivalence point.} All the acid is now \ce{HA-}. This
species is \textbf{amphoteric} --- it can lose another proton or gain
one back --- and it sits at the first jump.

\textbf{Second stage.} $\ce{HA- + OH- -> A^2- + H2O}$ takes the second
proton. The buffer region is now \ce{HA-}/\ce{A^2-}, and its half-point
gives pH $=$ p$K_{a2}$.

\textbf{Second equivalence point.} Everything is \ce{A^2-}, a
reasonably strong base, so the pH is well above 7.

\textbf{Two features to point out on a graph.} The second equivalence
volume is exactly \emph{twice} the first, because the second proton
needs as much base as the first. And the first jump is usually smaller
and less sharp than the second, because $K_{a1}$ and $K_{a2}$ are
typically only a few orders of magnitude apart --- the two stages
overlap slightly.

\textbf{What you will be asked.} Identify the equivalence points; say
which species predominates in each region; explain why the second
equivalence pH is higher. What you will \emph{not} be asked is to
compute the concentration of each species --- the exclusion note above
says so explicitly.
\end{workedexample}

\begin{yourturn}[YOUR TURN 14 --- four questions]
\begin{enumerate}[leftmargin=*,itemsep=0.5em]
  \item How many equivalence points does a diprotic acid curve have?
        \workspace[1.4cm]
  \item If the first is at \SI{20.0}{\milli\litre}, where is the
        second? \workspace[1.4cm]
  \item Which species predominates between the two equivalence points?
        \workspace[1.6cm]
  \item What does the pH equal at the first half-equivalence point?
        \workspace[1.6cm]
\end{enumerate}
\selfcheck{(a) two \quad (b) \SI{40.0}{\milli\litre} \quad (c)
\ce{HA-} \quad (d) p$K_{a1}$}
\keyonly{\begin{apcanswer}(a) \textbf{Two} --- one per ionisable
proton. (b) \textbf{\SI{40.0}{\milli\litre}}: the second proton
requires the same amount of base as the first, so the second
equivalence comes at twice the volume. (c) \textbf{\ce{HA-}}, the
amphoteric intermediate --- the first proton is gone and the second has
not yet been removed. (d) \textbf{p$K_{a1}$}, by the same argument as
Ladder 12 applied to the first ionisation.\end{apcanswer}}
\end{yourturn}

% =====================================================================
\section*{\sffamily Mixed set \textbullet\ ten questions, no labels}

\begin{apcnote}[Why this section exists]
Every question so far arrived with its ladder attached. These ten do
not --- and in this chapter the first decision is almost always
\emph{which region am I in}: excess strong reagent, buffer, or
equivalence. Build the ICF table and the region names itself.
\end{apcnote}

\begin{yourturn}[MIXED SET --- first five]
\begin{enumerate}[leftmargin=*,itemsep=0.5em,label=(\alph*)]
  \item A buffer is \SI{0.20}{\Molar} \ce{HA} and \SI{0.10}{\Molar}
        \ce{A-}, p$K_{a} = 4.74$. Find the pH. \workspace[1.8cm]
  \item Mix \SI{40.0}{\milli\litre} of \SI{0.100}{\Molar} \ce{HCl}
        with \SI{20.0}{\milli\litre} of \SI{0.100}{\Molar} \ce{NaOH}.
        Find the pH. \workspace[2.2cm]
  \item Which is a buffer: \ce{NH3}/\ce{NH4Cl} or
        \ce{NaOH}/\ce{NaCl}? \workspace[1.6cm]
  \item Equivalence is at \SI{36.0}{\milli\litre} and the pH at
        \SI{18.0}{\milli\litre} is 3.80. Find $K_{a}$.
        \workspace[2.0cm]
  \item Two buffers have the same ratio; one is five times more
        concentrated. Compare their pH and capacity.
        \workspace[1.8cm]
\end{enumerate}
\selfcheck{(a) 4.44 \quad (b) 1.48 \quad (c) \ce{NH3}/\ce{NH4Cl} \quad
(d) \num{1.6e-4} \quad (e) same pH, 5x capacity}
\keyonly{\begin{apcanswer}\textbf{(a)} pH $= 4.74 + \log(0.10/0.20) =
4.74 - 0.30 = \mathbf{4.44}$. \textbf{(b)} $n(\ce{H+}) =
\num{4.00e-3}$, $n(\ce{OH-}) = \num{2.00e-3}$, excess
$= \num{2.00e-3}$ mol in \SI{0.0600}{\litre}, so $[\ce{H+}] =
\num{3.33e-2}$ and pH $= \mathbf{1.48}$. \textbf{(c)}
\textbf{\ce{NH3}/\ce{NH4Cl}} --- a conjugate pair.
\ce{NaOH}/\ce{NaCl} is a strong base with an inert salt. \textbf{(d)}
\SI{18.0}{\milli\litre} is half-equivalence, so p$K_{a} = 3.80$ and
$K_{a} = 10^{-3.80} = \mathbf{\num{1.6e-4}}$. \textbf{(e)}
\textbf{Same pH} (ratio unchanged), \textbf{five times the capacity}
(five times as many moles of each component to consume added acid or
base).\end{apcanswer}}
\end{yourturn}

\begin{yourturn}[MIXED SET --- last five]
\begin{enumerate}[leftmargin=*,itemsep=0.5em,label=(\alph*)]
  \item \SI{50.0}{\milli\litre} of \SI{0.100}{\Molar} \ce{HA}
        (p$K_{a} = 4.74$) plus \SI{30.0}{\milli\litre} of
        \SI{0.100}{\Molar} \ce{NaOH}: find the pH.
        \workspace[2.2cm]
  \item The same acid plus \SI{50.0}{\milli\litre} of base: name the
        region and say whether the pH is above or below 7.
        \workspace[1.8cm]
  \item Which indicator would you use for that titration?
        \workspace[1.6cm]
  \item A weak base is titrated with strong acid. Is the equivalence pH
        above or below 7? \workspace[1.6cm]
  \item A diprotic acid's first equivalence is at
        \SI{15.0}{\milli\litre}. Where is the second, and which species
        dominates between them? \workspace[2.0cm]
\end{enumerate}
\selfcheck{(a) 4.92 \quad (b) equivalence, above 7 \quad (c)
phenolphthalein \quad (d) below 7 \quad (e)
\SI{30.0}{\milli\litre}; \ce{HA-}}
\keyonly{\begin{apcanswer}\textbf{(a)} ICF: \ce{HA} $= 5.00 - 3.00 =
\num{2.00e-3}$, \ce{A-} $= \num{3.00e-3}$; pH $= 4.74 +
\log(3.00/2.00) = 4.74 + 0.18 = \mathbf{4.92}$. \textbf{(b)} The
\textbf{equivalence point} (\SI{5.00e-3}{\mole} of base exactly
consumes the acid), and the pH is \textbf{above 7} because only \ce{A-},
a weak base, remains. \textbf{(c)} \textbf{Phenolphthalein} --- its
8.3--10.0 range brackets the basic equivalence pH. \textbf{(d)}
\textbf{Below 7}: the conjugate acid \ce{BH+} is left in the flask.
\textbf{(e)} \textbf{\SI{30.0}{\milli\litre}} --- twice the first ---
and \textbf{\ce{HA-}}, the amphoteric intermediate, predominates
between them.\end{apcanswer}}
\end{yourturn}

% =====================================================================
\section*{\sffamily Mastery tracker}

Tick a row only if \textbf{all four} YOUR TURN questions were right on
the first attempt. Two rows carry CED exclusions --- see the scope notes
in Ladders 6 and 14 --- but both exclusions are narrow, and everything
else here is fully assessed.

\renewcommand{\arraystretch}{1.25}
\noindent\begin{tabular}{@{}c p{5.7cm} c p{4.9cm}@{}}
\toprule
\textbf{First try?} & \textbf{Skill} & \textbf{Ladder} &
  \textbf{If not, re-read\ldots}\\
\midrule
$\square$ & the common-ion effect & 1 & adding a product shifts left \\
$\square$ & what a buffer is & 2 & both halves of one pair \\
$\square$ & Henderson--Hasselbalch & 3 & base on top, use p$K_{a}$ \\
$\square$ & designing a buffer & 4 & pick p$K_{a}$ near the target \\
$\square$ & buffer capacity & 5 & ratio sets pH, amount sets capacity \\
$\square$ & strong $+$ strong & 6 & moles first, total volume \\
$\square$ & the ICF stoichiometry step & 7 & F, not E \\
$\square$ & the three regions & 8 & the F row picks the method \\
$\square$ & the equivalence point & 9 & not pH 7 --- use $K_{b}$ \\
$\square$ & weak base $+$ strong acid & 10 & convert $K_{b}$ to $K_{a}$ \\
$\square$ & reading titration curves & 11 & same volume, different shape \\
$\square$ & the half-equivalence point & 12 & pH $=$ p$K_{a}$ \\
$\square$ & choosing an indicator & 13 & bracket the equivalence pH \\
$\square$ & polyprotic curves & 14 & shapes yes, computations no \\
\midrule
$\square$ & \textbf{mixed set} (8 of 10) & --- & whichever it exposed \\
\bottomrule
\end{tabular}

\begin{apcnote}[Scoring yourself honestly]
14/14 and the mixed set: Unit 8 is complete, and it is the largest
single block of marks in the course after Unit 3.

10--13: look at \emph{where}. Ladders 1--5 are buffers and repair
quickly. Ladders 6--10 are the two-step method, and nearly every miss
there is one of two things: forgetting the \textbf{total} volume, or
trying to do stoichiometry and equilibrium in the same table. Both are
mechanical and both are fixable in an afternoon.

9 or fewer: do not practise more problems yet. Go back to Ladder 7 and
build the ICF table for five different volumes of the same titration,
stopping each time at the F row without computing a pH. The skill this
chapter actually tests is \emph{recognising which of three situations
you are in}; the arithmetic afterwards is Chapter 14's and you already
have it.

\textbf{Where this sits.} Chapters 14 and 15 together are Unit 8 ---
11--15\% of the exam. With this chapter done, the self-study series is
complete: every chapter of Zumdahl that the CED assesses now has one.
\end{apcnote}

\end{document}
